\section{Interpretation and Scope}\label{sec:interpretation}

The mechanism identified above is not that firms make a small adjustment to quality and thereby marginally soften price competition. It is a discrete change in the relevant continuation game. A sufficiently large quality investment changes which pricing regime is played: the market moves from an interior shared-market equilibrium to an exclusion equilibrium. Once upstream actions can move the downstream game across an active-set boundary, a first-order condition computed on one smooth branch need not identify a global equilibrium.

The timing in the original model has a natural economic reading. Variety or physical position is the slowest choice, quality is a medium-run commitment, and price adjusts fastest. \citet{Economides1989} motivates the quality index with examples such as advertising, computer performance, and warranty provision. We retain that interpretation only as motivation for the ordering of decisions. The present paper does not claim that any specific observed industry follows the exact linear Hotelling environment or that actual quality competition necessarily produces exclusion.

The model nonetheless suggests a qualitative empirical signature. If quality investment becomes sufficiently attractive relative to the intensity of horizontal differentiation, adjustment need not be smooth. Quality dispersion, market-share concentration, and price dispersion can move together when one firm's quality choice crosses the threshold that changes the downstream pricing regime. This is a model implication rather than an empirical finding.

The scope is intentionally narrow. We do not add heterogeneous quality tastes, alternative transportation technologies, entry, spillovers, or mixed strategies. The point of the exercise is to reassess the equilibrium logic of the original game before introducing new mechanisms. In particular, the paper establishes existence of specific asymmetric quality equilibria but does not characterize the entire pure quality-equilibrium correspondence.
